% Vietnamese translation for OI Public Library
% Source-PDF: 动态规划 DYNAMIC PROGRAMMING/树形DP选讲 - 顾逸宏.pdf
\documentclass[11pt]{article}
\usepackage{oipl-vn}

\title{Chuyên đề quy hoạch động trên cây}
\author{Cố Dật Hoành}
\date{16 tháng 8, 2016}

\begin{document}
\maketitle
\sourcepath{Dynamic Programming / Tree DP selected topics - Gu Yihong}

\begin{abstract}
Bản gốc là một bộ slide Beamer: 142 trang PDF chủ yếu đến từ các lớp phủ
từng bước của khoảng 40 slide nội dung. Bản tiếng Việt này chuyển slide thành
ghi chú bài giảng tiếng Việt, giữ các trạng thái, công thức chuyển và ví dụ
thuật toán chính; những hình minh họa đơn giản được diễn giải trong phần văn
bản.
\end{abstract}

\tableofcontents

\section{Cách nắm chắc quy hoạch động}

Cốt lõi của quy hoạch động là trạng thái. Một trạng thái gồm hai phần:
cách biểu diễn và cách chuyển.

Năng lực quan trọng nhất khi làm quy hoạch động là phân loại trường hợp và
quy nạp. Vì vậy nên nắm chắc các dạng cơ bản: quy hoạch động thông thường,
quy hoạch động trên cây, quy hoạch động chữ số, quy hoạch động nén trạng thái
và các kỹ thuật tối ưu hóa quy hoạch động. Phần còn lại đến từ luyện tập nhiều
bài, chẳng hạn các bài trên Topcoder.

Trong tài liệu này dùng các ký hiệu sau:
\begin{itemize}
  \item \(x\): đỉnh hiện đang xét;
  \item \(c\): một con của \(x\);
  \item \(f\): cha của \(x\);
  \item \((x,y)\): đường đi từ \(x\) đến \(y\);
  \item \(w(x,y)\): độ dài đường đi từ \(x\) đến \(y\).
\end{itemize}

\section{Bài toán kinh điển: đường đi ngắn nhất lớn nhất}

Bài toán mẫu xuyên suốt các slide là: cho một đồ thị vô hướng \(G\), hãy tìm
giá trị lớn nhất của khoảng cách ngắn nhất giữa hai đỉnh bất kỳ. Trên cây,
đó chính là đường kính cây; trên đồ thị một chu trình hoặc cactus, cần xử lý
thêm phần chu trình.

\subsection{Trường hợp 1: đồ thị thông thường}

Cách thô là liệt kê đỉnh xuất phát \(i\), chạy SPFA hoặc một thuật toán đường
đi ngắn nhất thích hợp, rồi lấy giá trị lớn nhất. Có thể dùng Floyd với độ
phức tạp \(O(n^3)\) nếu miền dữ liệu nhỏ và đồ thị đủ dày.

\subsection{Trường hợp 2: cây}

\subsubsection{Cách tham lam}

Chọn tùy ý một đỉnh \(x\), tìm đỉnh \(y\) xa \(x\) nhất trong cây. Sau đó tìm
đỉnh \(z\) xa \(y\) nhất. Đường đi từ \(y\) đến \(z\) là một đường kính của cây.

Ý tưởng chứng minh là phản chứng. Nếu tồn tại một đường đi \((y',z')\) dài hơn
\((y,z)\), thì các đỉnh \(y,z,y',z'\) đều là lá. Xét theo việc hai đường đi
\((y,z)\) và \((y',z')\) có giao nhau hay không, luôn suy ra mâu thuẫn với cách
chọn \(y\) hoặc \(z\) là đỉnh xa nhất.

\subsubsection{Cách quy hoạch động}

Gốc hóa cây tại đỉnh \(1\). Gọi \(L[x]\) là độ dài đường đi dài nhất nằm hoàn
toàn trong cây con gốc \(x\). Khi đó \(L[1]\) là đáp án.

Để tính \(L[x]\), phân loại đường đi dài nhất trong cây con của \(x\):
\begin{itemize}
  \item đường đi không qua \(x\): lấy giá trị lớn nhất trong các \(L[c]\);
  \item đường đi qua \(x\): dùng hai nhánh đi xuống tốt nhất từ các con của
  \(x\).
\end{itemize}

Gọi \(down[x]\) là độ dài đường đi dài nhất bắt đầu tại \(x\) và đi xuống trong
cây con. Khi đó
\[
  down[x] = \max_c \{ down[c] + w(x,c) \}.
\]
Đường đi tốt nhất qua \(x\) là tổng của giá trị lớn nhất và lớn nhì trong các
\(down[c] + w(x,c)\), trong đó giá trị lớn nhì có thể khởi tạo bằng \(0\).

\subsection{Mở rộng: \texorpdfstring{\(down\)}{down} và \texorpdfstring{\(up\)}{up}}

Bài toán mở rộng: \(G\) là cây, đồng thời cần biết với mỗi đỉnh \(x\), đường đi
dài nhất bắt đầu từ \(x\) là bao nhiêu.

Phân loại bước đầu tiên:
\begin{itemize}
  \item đi xuống: đóng góp là \(down[x]\);
  \item đi lên: gọi \(up[x]\) là đường đi dài nhất bắt đầu tại \(x\) với bước
  đầu tiên đi lên cha.
\end{itemize}

Để tính \(up[x]\), giả sử bước đầu tiên đi từ \(x\) lên \(f\). Bước tiếp theo có
hai khả năng:
\begin{itemize}
  \item tiếp tục đi lên: dùng \(up[f]\);
  \item đi xuống một nhánh khác của \(f\): nếu \(down[f]\) không được chuyển
  từ chính \(x\), dùng \(down[f]\); nếu không, dùng giá trị tốt thứ hai
  \(down2[f]\).
\end{itemize}
Vì vậy
\[
  up[x] = w(x,f) + \max\{up[f],\; best\_down\_from\_f\_excluding\_x\}.
\]
Cần lưu thêm \(down2[x]\), giá trị lớn nhì của \(down[c]+w(x,c)\), và
\(down1v[x]\), đứa con tạo ra \(down[x]\).

\subsection{Trường hợp 3: đồ thị một chu trình}

Giả sử \(G\) chỉ có đúng một chu trình. Cần chú ý rằng bài toán vẫn hỏi giá trị
lớn nhất của khoảng cách ngắn nhất. Trước hết chạy DFS hoặc BFS để tìm chu
trình, ghi các đỉnh trên chu trình là \(cir[i]\).

Phân loại đường đi:
\begin{itemize}
  \item không đi qua chu trình: xử lý như trên cây và lấy lớn nhất trong các
  \(L[cir[i]]\);
  \item đi qua chu trình: chọn hai đỉnh \(x,y\) trên chu trình, cộng hai nhánh
  cây treo và đoạn ngắn hơn trên chu trình.
\end{itemize}

Cách thô liệt kê hai đỉnh trên chu trình:
\[
  down[x] + down[y] + \min\{w_1(x,y), w_2(x,y)\},
\]
với \(w_1,w_2\) là hai khoảng cách theo hai chiều trên chu trình. Độ phức tạp là
\(O(k^2)\), trong đó \(k\) là độ dài chu trình.

Cách tối ưu: cắt chu trình thành một dãy và nhân đôi dãy đó. Chọn một hướng cố
định. Khi liệt kê đỉnh \(x\), các đỉnh \(y\) mà đi theo hướng này đến \(x\) là
đoạn ngắn hơn tạo thành một khoảng \([l_x,r_x]\). Hai biên \(l_x,r_x\) đều đơn
điệu không giảm, nên dùng hàng đợi đơn điệu để duy trì cực trị. Độ phức tạp còn
\(O(k)\).

\subsection{Trường hợp 4: cactus}

Ở đây \(G\) là cactus, nghĩa là hai chu trình bất kỳ có nhiều nhất một đỉnh
chung. Chạy DFS một lần để lấy thứ tự DFS. Với mỗi đỉnh \(x\), xét từng cạnh
\((x,y)\):
\begin{itemize}
  \item nếu là cạnh cây, cập nhật như quy hoạch động trên cây;
  \item nếu không phải cạnh cây, cạnh này thuộc một chu trình cần xử lý riêng.
\end{itemize}

Trên mỗi chu trình, làm từ thứ tự DFS lớn về nhỏ. Thứ tự tính đáp án và cập
nhật \(down\) của một chu trình liên quan đến đỉnh có thứ tự DFS nhỏ nhất trong
chu trình đó. Các cạnh nằm trên chu trình không trực tiếp cập nhật \(down\).
Trong một chu trình, phần có thể ảnh hưởng đến đáp án phía sau chỉ là đỉnh có
thứ tự DFS nhỏ nhất.

Đoạn mã lõi trong slide được giữ nguyên về ý tưởng và tên biến:

\begin{verbatim}
for (int k = c[x]; ~k; k = nxt[k]) {
    int y = g[k];
    if (fa[y] == x && low[y] > idx[x]) {
        ans = max(ans, down[y] + 1 + down[x]);
        down[x] = max(down[x], down[y] + 1);
    }
    if (fa[y] != x && idx[x] < idx[y]) {
        N = 0;
        for (int i = y; i != fa[x]; i = fa[i])
            cir[++N] = i;
        updans();
        for (int i = 1; i < N; ++i)
            ckmax(down[x], down[cir[i]] + min(i, N - i));
    }
}
\end{verbatim}

\section{Kiến thức chuẩn bị}

Các slide liệt kê một số nền tảng nên có trước khi học sâu phần này:
\begin{itemize}
  \item tổ hợp và các cách khử trùng lặp đơn giản;
  \item gộp theo heuristic, thường là gộp nhỏ vào lớn;
  \item băm chuỗi;
  \item xác suất, biến ngẫu nhiên rời rạc và kỳ vọng.
\end{itemize}

\section{Ví dụ kinh điển}

\subsection{SPOJ MTREE}

Cho một cây, tính tổng tích trọng số cạnh trên mọi đường đi, lấy modulo. Giới
hạn \(N \le 10^5\). Đây là dạng quy hoạch động cộng dồn đóng góp đường đi qua
mỗi đỉnh hoặc mỗi cạnh.

\subsection{SPOJ GS}

Cho một cây. Tại mỗi đỉnh, xác suất đi sang từng đỉnh kề đã được cho trước.
Hỏi kỳ vọng số bước đi từ một đỉnh đến một đỉnh khác. Giới hạn
\(N \le 10^5\). Dạng này thường dẫn đến hệ phương trình hoặc hai lượt DP trên
cây để biểu diễn kỳ vọng theo cha/con.

\subsection{CodeChef TAPAIR}

Cho một đồ thị vô hướng liên thông \(N\) đỉnh. Hỏi có bao nhiêu cách xóa hai
cạnh để đồ thị trở nên không liên thông. Giới hạn \(N,M \le 10^5\). Bài liên
quan đến cầu, cây DFS và việc đếm cặp cạnh theo cấu trúc khối.

\subsection{SPOJ TREECST}

Xóa một cạnh của cây rồi thêm một cạnh khác để vẫn được một cây, sao cho đường
kính của cây mới nhỏ nhất; cần in cả phương án. Giới hạn \(N \le 3\cdot 10^5\).
Các đại lượng đường kính, tâm cây và giá trị \(down/up\) thường là phần lõi.

\subsection{SPOJ TWOPATHS}

Trên một cây, tìm hai đường đi không giao nhau nghiêm ngặt sao cho tích độ dài
của chúng là lớn nhất, rồi in tích đó. Giới hạn \(N \le 10^5\). Cần kết hợp
đường kính trong cây con với thông tin tốt nhất bên ngoài hoặc theo cạnh cắt.

\subsection{Codeforces 202 Div1 E}

Cho một cây \(n\) đỉnh, trong đó có \(m\) đỉnh là tu viện. Mỗi tu viện lập danh
sách các tu viện xa nó nhất và chuẩn bị đi thăm tất cả các nơi đó. Một đỉnh
không có tu viện có thể bị xóa; nếu điều đó khiến một tu viện không còn nơi nào
để đi thăm, giá trị vui thích tăng thêm \(1\). Hỏi giá trị vui thích lớn nhất và
số phương án đạt giá trị đó. Giới hạn \(m \le n \le 10^5\).

\subsection{NOI2012 Day2: Công viên lạc lối}

Cho đồ thị vô hướng liên thông \(n\) đỉnh, \(m\) cạnh. Ban đầu đứng ngẫu nhiên
ở một đỉnh; mỗi bước chọn đều trong các đỉnh kề chưa đi qua. Hỏi kỳ vọng số
bước di chuyển. Đề bảo đảm \(m=n-1\) hoặc \(m=n\); nếu \(m=n\), đồ thị có đúng
một chu trình và chu trình có kích thước không quá \(30\), với \(n\le 10^5\).

\section{Bài suy nghĩ}

\subsection{FoxConnection}

Cho một cây \(n\) đỉnh; mỗi đỉnh có thể có một con fox. Cần di chuyển các con
fox sao cho mỗi đỉnh có nhiều nhất một con, tập các đỉnh có fox liên thông, và
tổng khoảng cách di chuyển là nhỏ nhất. Giới hạn \(n \le 50\).

\subsection{InducedSubgraphs}

Cho một cây \(n\) đỉnh. Cần đánh nhãn lại các đỉnh sao cho với mỗi tập liên tiếp
\[
  \{1,2,\ldots,k\},\ \{2,3,\ldots,k+1\},\ \ldots,\ 
  \{n-k+1,n-k+2,\ldots,n\},
\]
cây con cảm sinh bởi tập đó đều có đúng \(k\) đỉnh. Giới hạn \(n\le 50\).

Slide tách hai trường hợp \(2k\le n\) và \(2k\ge n\). Với \(2k\le n\), cấu trúc
bắt buộc là một chuỗi dài ở giữa, hai đầu treo các cây con kích thước \(k\):
\[
  [STA] - (k+1) - (k+2) - \cdots - (n-k) - [STB].
\]
Ta liệt kê chuỗi giữa, tìm hai phần \(STA\) và \(STB\). Việc tô màu trong
\(STA\) tương đương yêu cầu mọi con đều có nhãn nhỏ hơn cha, từ đó đếm số cách
tô màu. Slide ký hiệu \(f[x][p]\) là số cách tô màu cây con gốc \(x\) khi cha
của \(x\) mang màu \(p\); phần \(STB\) xử lý đối xứng.

Với trường hợp \(2k\ge n\), các tập liên tiếp độ dài \(k\) chồng lấn lên nhau.
Phần giao bắt buộc là một khối liên thông lớn có kích thước \(2k-n\). Các cây
con treo ra khỏi khối này phải được phân về một trong hai phía nhãn: hoặc
thuộc đoạn \([1,v]\), hoặc thuộc đoạn \([n-v+1,n]\), tùy phía của đoạn liên
tiếp đang xét. Sau khi chọn khối trung tâm, ta dùng một quy hoạch động kiểu
ba lô hai chiều để phân phối kích thước các cây con treo cho hai phía. Phần
đếm nhãn trong khối trung tâm cần chia cho số cách quay vòng tương đương mà
slide ghi dưới dạng hiệu chỉnh
\[
  \frac{(2k-n)!}{2k-n}.
\]

\subsection{EagleInZoo}

Có một cây và \(k\) con đại bàng sẽ lần lượt bay đến một đỉnh để nghỉ. Mỗi con
đại bàng bắt đầu từ gốc \(0\). Nếu đỉnh hiện tại còn trống thì nó dừng lại tại
đó. Nếu đỉnh hiện tại đã có đại bàng, nó chọn ngẫu nhiên đều một đứa con để
bay tiếp; nếu không có con thì nó chỉ có thể bay đi khỏi cây. Cần tính xác
suất con đại bàng thứ \(k\) có thể dừng lại trên cây. Giới hạn trong slide là
\[
  n\le 50,\qquad k\le 100.
\]

Phân loại theo đỉnh dừng của con đại bàng thứ \(k\). Gọi \(P(i)\) là xác suất
con thứ \(k\) dừng ở đỉnh \(i\), khi đó
\[
  ans=\sum_{i=1}^{n} P(i).
\]
Đặt \(f[i][j]\) là xác suất sau \(i\) lượt, đỉnh \(j\) vẫn trống nhưng cha của
\(j\) đã không trống. Gọi \(\mathrm{reach}[j]\) là xác suất đi từ gốc đến
\(j\) theo các lựa chọn ngẫu nhiên trên cây. Khi đó đóng góp của đỉnh \(j\)
cho lượt thứ \(K\) là
\[
  f[K-1][j]\cdot \mathrm{reach}[j],
\]
nên
\[
  ans=\sum_{j=1}^{n} f[K-1][j]\cdot \mathrm{reach}[j].
\]
Các chuyển trạng thái chỉ cần xét việc một lượt mới có đi qua nhánh chứa
\(j\) hay không, nên slide kết luận chuyển ``rất đơn giản'' và tổng độ phức
tạp là \(O(NK)\).

\section*{Lời cảm ơn}

Slide cuối ghi lời cảm ơn và địa chỉ liên hệ của tác giả:
\texttt{sbullet@163.com}.

\end{document}
